% GENERATED FILE. Edit canonical lessons, then run npm run generate:books.
% canonical-source-hash: fnv1a64:e2cdda09df99fe9d
% canonical-lessons: GU-C428-devalu, GU-C428-lenu, GU-C428-uchinu, GU-C428-sahukar, GU-C428-thagvu

\chapter{Debts and Moneylenders}
\label{ch:gu-debts-and-moneylenders}
\hlchaptermodality{428}

\begin{chapteropening}
\textbf{By the end of this chapter:} \emph{I can talk about owing and lending money.}

Complete the last lesson of chapter 428: I can talk about owing and lending money.
\end{chapteropening}

\section[\texorpdfstring{\gu{દેવાળું} (devāḷũ)}{દેવાળું (devāḷũ)}]{\gu{દેવાળું} (devāḷũ) — bankruptcy}

\label{lesson:GU-C428-devalu}



\begin{quote}
Before the new one: say the Gujarati for a fixed term, then the Gujarati for to allot.
\end{quote}

\subsection*{You'll want to know: \gu{દેવાળું}}
\textbf{\gu{દેવાળું}} — \emph{devāḷũ} — \textquotedblleft{}bankruptcy\textquotedblright{}.

\begin{quote}
\gu{ધંધામાં} \gu{ખોટ} \gu{ગઈ} \gu{અને} \gu{વેપારીનું} \gu{દેવાળું} \gu{નીકળ્યું}.
\end{quote}

\begin{quote}
\emph{dhandhāmā̃ khoṭ gaī ane vepārīnũ devāḷũ nīkaḷyũ.}
\end{quote}

\begin{quote}
The business made a loss and the merchant went bankrupt.
\end{quote}

A neuter noun.

\begin{cousinweb}
\textbf{\gu{દેવાળું}} (\emph{devāḷũ}) CDIAL lists Hindi \emph{diwālā} and Punjabi \emph{davālā}, \textquotedblleft{}bankruptcy\textquotedblright{}, under Sanskrit \emph{dīpāvali}, \textquotedblleft{}a row of lamps\textquotedblright{} (CDIAL 6358), the word behind \gu{દિવાળી}; it gives no Gujarati form for \gu{દેવાળું}, and how lamps came to mean ruin is not settled, so the link is hedged.
\end{cousinweb}

\begin{grammarlens}[title={What you've built}]
One more word for getting around.
\end{grammarlens}

\subsection*{Your turn}
\begin{itemize}
\raggedright
  \item \emph{Say it:} \emph{devāḷũ}
  \item \emph{Say it:} \emph{devāḷũ}, once more
  \item \emph{Say it:} \emph{dhandhāmā̃ khoṭ gaī ane vepārīnũ devāḷũ nīkaḷyũ.} — The business made a loss and the merchant went bankrupt.
  \item \emph{Recall:} say the Gujarati for a fixed term, then the Gujarati for to allot, then say \emph{devāḷũ} again
\end{itemize}

\begin{tcolorbox}[breakable,colback=teal!4,colframe=teal!35!black,title={Before you move on}]
What does \gu{દેવાળું} mean? (\textquotedblleft{}Bankruptcy\textquotedblright{}.) Say it once more.
\end{tcolorbox}
\section[\texorpdfstring{\gu{લેણું} (leṇũ)}{લેણું (leṇũ)}]{\gu{લેણું} (leṇũ) — money owed to you; a sum due}

\label{lesson:GU-C428-lenu}



\begin{quote}
Before the new one: say the Gujarati for to allot, then the Gujarati for bankruptcy.
\end{quote}

\subsection*{You'll want to know: \gu{લેણું}}
\textbf{\gu{લેણું}} — \emph{leṇũ} — \textquotedblleft{}money owed to you; a sum due\textquotedblright{}.

\begin{quote}
\gu{તેની} \gu{પાસે} \gu{મારું} \gu{સો} \gu{રૂપિયાનું} \gu{લેણું} \gu{છે}.
\end{quote}

\begin{quote}
\emph{tenī pāse mārũ so rūpiyānũ leṇũ che.}
\end{quote}

\begin{quote}
He owes me a hundred rupees.
\end{quote}

A neuter noun.

\begin{cousinweb}
\textbf{\gu{લેણું}} (\emph{leṇũ}) is built on \gu{લેવું}, to take, a tadbhava of Sanskrit \emph{labhatē}, \textquotedblleft{}catches, takes\textquotedblright{} (CDIAL 10948): what you still have to take back from someone.
\end{cousinweb}

\begin{grammarlens}[title={What you've built}]
One more word for getting around.
\end{grammarlens}

\subsection*{Your turn}
\begin{itemize}
\raggedright
  \item \emph{Say it:} \emph{leṇũ}
  \item \emph{Say it:} \emph{leṇũ}, once more
  \item \emph{Say it:} \emph{tenī pāse mārũ so rūpiyānũ leṇũ che.} — He owes me a hundred rupees.
  \item \emph{Recall:} say the Gujarati for to allot, then the Gujarati for bankruptcy, then say \emph{leṇũ} again
\end{itemize}

\begin{tcolorbox}[breakable,colback=teal!4,colframe=teal!35!black,title={Before you move on}]
What does \gu{લેણું} mean? (\textquotedblleft{}Money owed to you; a sum due\textquotedblright{}.) Say it once more.
\end{tcolorbox}
\section[\texorpdfstring{\gu{ઉછીનું} (uchīnũ)}{ઉછીનું (uchīnũ)}]{\gu{ઉછીનું} (uchīnũ) — borrowed for a while}

\label{lesson:GU-C428-uchinu}



\begin{quote}
Before the new one: say the Gujarati for bankruptcy, then the Gujarati for money owed to you.
\end{quote}

\subsection*{You'll want to know: \gu{ઉછીનું}}
\textbf{\gu{ઉછીનું}} — \emph{uchīnũ} — \textquotedblleft{}borrowed for a while\textquotedblright{}.

\begin{quote}
\gu{મેં} \gu{મિત્ર} \gu{પાસેથી} \gu{સો} \gu{રૂપિયા} \gu{ઉછીના} \gu{લીધા}.
\end{quote}

\begin{quote}
\emph{mẽ mitra pāsethī so rūpiyā uchīnā līdhā.}
\end{quote}

\begin{quote}
I borrowed a hundred rupees from a friend.
\end{quote}

\begin{cousinweb}
\textbf{\gu{ઉછીનું}} (\emph{uchīnũ}) goes back to \emph{*ut-śrita-}, \textquotedblleft{}raised\textquotedblright{}, a form CDIAL reconstructs beside Sanskrit \emph{ucchrita}; CDIAL lists the Gujarati word, \textquotedblleft{}borrowed\textquotedblright{} (CDIAL 1862).
\end{cousinweb}

\begin{grammarlens}[title={What you've built}]
One more word for getting around.
\end{grammarlens}

\subsection*{Your turn}
\begin{itemize}
\raggedright
  \item \emph{Say it:} \emph{uchīnũ}
  \item \emph{Say it:} \emph{uchīnũ}, once more
  \item \emph{Say it:} \emph{mẽ mitra pāsethī so rūpiyā uchīnā līdhā.} — I borrowed a hundred rupees from a friend.
  \item \emph{Recall:} say the Gujarati for bankruptcy, then the Gujarati for money owed to you, then say \emph{uchīnũ} again
\end{itemize}

\begin{tcolorbox}[breakable,colback=teal!4,colframe=teal!35!black,title={Before you move on}]
What does \gu{ઉછીનું} mean? (\textquotedblleft{}Borrowed for a while\textquotedblright{}.) Say it once more.
\end{tcolorbox}
\section[\texorpdfstring{\gu{શાહુકાર} (śāhukār)}{શાહુકાર (śāhukār)}]{\gu{શાહુકાર} (śāhukār) — a moneylender}

\label{lesson:GU-C428-sahukar}



\begin{quote}
Before the new one: say the Gujarati for money owed to you, then the Gujarati for borrowed for a while.
\end{quote}

\subsection*{You'll want to know: \gu{શાહુકાર}}
\textbf{\gu{શાહુકાર}} — \emph{śāhukār} — \textquotedblleft{}a moneylender\textquotedblright{}.

\begin{quote}
\gu{ગામનો} \gu{શાહુકાર} \gu{ઘણું} \gu{વ્યાજ} \gu{લે} \gu{છે}.
\end{quote}

\begin{quote}
\emph{gāmno śāhukār ghaṇũ vyāj le che.}
\end{quote}

\begin{quote}
The village moneylender charges a lot of interest.
\end{quote}

A masculine noun.

\begin{cousinweb}
\textbf{\gu{શાહુકાર}} (\emph{śāhukār}) is built on \gu{શાહુ}, a banker, which goes back to Sanskrit \emph{sādhu}, \textquotedblleft{}good, honest\textquotedblright{} (CDIAL 13337), with -\gu{કાર}, \textquotedblleft{}one who does\textquotedblright{}.
\end{cousinweb}

\begin{grammarlens}[title={What you've built}]
One more word for getting around.
\end{grammarlens}

\subsection*{Your turn}
\begin{itemize}
\raggedright
  \item \emph{Say it:} \emph{śāhukār}
  \item \emph{Say it:} \emph{śāhukār}, once more
  \item \emph{Say it:} \emph{gāmno śāhukār ghaṇũ vyāj le che.} — The village moneylender charges a lot of interest.
  \item \emph{Recall:} say the Gujarati for money owed to you, then the Gujarati for borrowed for a while, then say \emph{śāhukār} again
\end{itemize}

\begin{tcolorbox}[breakable,colback=teal!4,colframe=teal!35!black,title={Before you move on}]
What does \gu{શાહુકાર} mean? (\textquotedblleft{}A moneylender\textquotedblright{}.) Say it once more.
\end{tcolorbox}
\section[\texorpdfstring{\gu{ઠગવું} (ṭhagvũ)}{ઠગવું (ṭhagvũ)}]{\gu{ઠગવું} (ṭhagvũ) — to swindle; to con}

\label{lesson:GU-C428-thagvu}



\begin{quote}
Before the new one: say the Gujarati for borrowed for a while, then the Gujarati for a moneylender.
\end{quote}

\subsection*{You'll want to know: \gu{ઠગવું}}
\textbf{\gu{ઠગવું}} — \emph{ṭhagvũ} — \textquotedblleft{}to swindle; to con\textquotedblright{}.

\begin{quote}
\gu{આ} \gu{દુકાનદાર} \gu{ગ્રાહકોને} \gu{ઠગે} \gu{છે}.
\end{quote}

\begin{quote}
\emph{ā dukāndār grāhkone ṭhage che.}
\end{quote}

\begin{quote}
This shopkeeper swindles his customers.
\end{quote}

Past, with \gu{તેણે} (he or she …): \gu{તેણે} \gu{વેપારીને} \gu{ઠગ્યો}, \emph{teṇe vepārīne ṭhagyo}.

\begin{cousinweb}
\textbf{\gu{ઠગવું}} (\emph{ṭhagvũ}) goes back to \emph{*ṭhagg}, \textquotedblleft{}to cheat\textquotedblright{}, a form CDIAL reconstructs; CDIAL lists the Gujarati verb, \textquotedblleft{}to cheat, to rob\textquotedblright{} (CDIAL 5489).
\end{cousinweb}

\begin{grammarlens}[title={What you've built}]
One more word for getting around.
\end{grammarlens}

\subsection*{Your turn}
\begin{itemize}
\raggedright
  \item \emph{Say it:} \emph{ṭhagvũ}
  \item \emph{Say it:} \emph{ṭhagvũ}, once more
  \item \emph{Say it:} \emph{ā dukāndār grāhkone ṭhage che.} — This shopkeeper swindles his customers.
  \item \emph{Recall:} say the Gujarati for borrowed for a while, then the Gujarati for a moneylender, then say \emph{ṭhagvũ} again
\end{itemize}

\begin{tcolorbox}[breakable,colback=teal!4,colframe=teal!35!black,title={Before you move on}]
What does \gu{ઠગવું} mean? (\textquotedblleft{}To swindle; to con\textquotedblright{}.) Say it once more.
\end{tcolorbox}
